\section{Method}
\label{sec:method}

\subsection{Study design}
% TODO: one teacher (E0), one student (E1, E2); every compression arm starts from E2
% (or E7 for E8, E0 for E9); all arms are compared at the same MEASURED budget
% (peak memory or latency as a fraction of E2), not parameter count.
\todo{study design}

\subsection{Compression arms}
% TODO: describe E3 (PTQ: W8, W4 via GPTQ/AWQ, W8A8), E4 (QAKD, GRACE-style),
% E5 (magnitude, activation, Taylor), E6 (output-similarity layer pruning), E7 (ours),
% E8 (E7 + 4-bit), E9 (teacher 4-bit).

\begin{table}[t]
  \centering
  \caption{Experiment ladder (\expid{E0}--\expid{E9}) and ablations
  (\expid{A1}--\expid{A6}). Source: \texttt{src/medcompress/registry.py}.}
  \label{tab:ladder}
  \begin{tabular}{llll}
    \toprule
    ID & Description & Starts from & Budgets \\
    \midrule
    \expid{E0} & Teacher, full precision & -- & -- \\
    \expid{E1} & Student, answer fine-tuning & -- & -- \\
    \expid{E2} & Student + offline KD & -- & reference \\
    \expid{E3} & PTQ (W8/W4/W8A8) & \expid{E2} & measured \\
    \expid{E4} & QAKD & -- & measured \\
    \expid{E5} & Generic structured pruning + recovery & \expid{E2} & B1--B3 \\
    \expid{E6} & Output-similarity layer pruning + recovery & \expid{E2} & B1--B3 \\
    \expid{E7} & Clinical pruning (ours) + recovery & \expid{E2} & B1--B3 \\
    \expid{E8} & \expid{E7} + 4-bit PTQ & \expid{E7} & B1--B3 \\
    \expid{E9} & Teacher 4-bit (suppl.) & \expid{E0} & measured \\
    \expid{A1}--\expid{A6} & Criterion ablations & \expid{E2} & B2 \\
    \bottomrule
  \end{tabular}
\end{table}

\subsection{Clinically weighted pruning criterion}
% TODO: units (layers, GQA-consistent head groups, MLP channel groups; projector kept),
% probes, Delta variants, weights, procedure. Keep consistent with
% src/medcompress/compression/pruning/criteria/clinical.py.
For each candidate unit $j$ we compute
\begin{equation}
  I_j = \sum_{c \in C} w_c \, \frac{1}{|X_{\mathrm{cal}}|} \sum_{x \in X_{\mathrm{cal}}} \Delta_j(c, x),
  \label{eq:importance}
\end{equation}
where $\Delta_j(c, x) = \log p_\theta(y_{c,x} \mid x, q_c) - \log p_{\theta \setminus j}(y_{c,x} \mid x, q_c)$
is the drop in log-probability of the correct yes/no answer to the probe $q_c$ when
unit $j$ is ablated, and $w_c = \min(1/\pi_c, \kappa)\, \tau(c)$ combines capped
inverse prevalence and a pre-registered clinical importance tier.
\todo{full definitions, KL variant, procedure}

\begin{figure}[t]
  \centering
  \fbox{\parbox{0.9\linewidth}{\centering\vspace{3em}
    Placeholder: probe $\rightarrow$ ablate $\rightarrow$ score $\rightarrow$ rank
    $\rightarrow$ remove to budget $\rightarrow$ recover $\rightarrow$ mask audit.\vspace{3em}}}
  \caption{Overview of the clinical pruning criterion (\expid{E7}).}
  \label{fig:method}
\end{figure}

\subsection{Shared recovery}
% TODO: identical LoRA + offline KD recovery for all pruning arms (E5, E6, E7,
% A1-A6): same data subset, steps and hyper-parameters; recovery config fingerprint.
\todo{recovery}

\subsection{Mask audit}
% TODO: lesion masked vs area-matched random region masked; lesion-mask effect.
\todo{mask audit}
